\section{Discussion}

\subsection{Key Takeaways}

\subsubsection{Takeaway 1: Real Compression, Zero-Shot}

SOA Fused Attention \textbf{reduces} attention parameters by \textbf{1.92$\times$} 
without requiring any calibration data. This is a significant practical 
advantage over A3 and AASVD.

\subsubsection{Takeaway 2: $W_O$ Should Stay}

Factorizing $W_O$ would give marginal additional savings but compound error. 
We keep it intact.

\subsubsection{Takeaway 3: Works Across Architectures}

The same \texttt{copy\_weights} handles DistilBERT, RoBERTa, GPT-2, and Qwen3.5.

\subsection{When to Use SOA Fused Attention}

\subsubsection{Recommended}
\begin{itemize}
\item \textbf{Zero-shot deployment} where calibration data is unavailable.
\item \textbf{Parameter-sensitive deployments} where every million counts.
\item \textbf{Attention-heavy models} (DistilBERT, BERT-family).
\item \textbf{Composing with GQA} for additional KV reduction.

\end{itemize}
\subsubsection{Not Recommended}
\begin{itemize}
\item \textbf{Long-context inference} -- attention is still $\mathcal{O}(L^2)$.
\item \textbf{Models with small attention share} (e.g., large FFN expansion).
\item \textbf{Very low rank} ($k < 64$) -- reconstruction error too high.

\end{itemize}
\subsection{Decision Flowchart}

\begin{lstlisting}
Have calibration data? --- Yes --- Consider A3 or AASVD
  | No
   
Want to reduce parameters? --- Yes --- Use SOA Fused
  | No
   
Consider attention-replacement methods
\end{lstlisting}

\subsection{Cost-Benefit Analysis}

For DistilBERT at $k = 128$:

\begin{table}[H]
\centering
\small
\begin{tabular}{llll}
\toprule
\textbf{Metric} & \textbf{Baseline} & \textbf{Ours} & \textbf{Change} \\
\midrule
Attention params & 2.36M & 1228800 & \textbf{1.92$\times$} \\
Total params & 66.95M & 65536000 & -2.1\% \\
Inference time & 1.90 s & 2.1 s & 1.1$\times$ \\
Accuracy & 90.71\% & 81.31\% & -9.40\% \\
\bottomrule
\end{tabular}
\end{table}

\subsection{Composition with Other Methods}

SOA Fused Attention composes with:

\begin{itemize}
\item \textbf{GQA}: further reduces KV projection size.
\item \textbf{KV quantization}: INT8 on the factorized projections.
\item \textbf{Token compression} (TCA): reduces KV length after projection.
\item \textbf{FFN compression} (A3/AASVD): compresses the other half of the block.

\end{itemize}
A combined pipeline could achieve \textbf{4-5$\times$ total compression} with the same 
accuracy recovery strategy.

\subsection{Broader Impact}

\subsubsection{Positive}
\begin{itemize}
\item Enables deployment on smaller hardware.
\item Zero-shot -- no privacy concerns from calibration data.
\item Simple to apply.

\end{itemize}
\subsubsection{Negative}
\begin{itemize}
\item Attention remains quadratic -- not a long-context solution.
\item $W_O$ preserved -- modest total reduction.

\end{itemize}
\subsection{Open Questions}

\begin{enumerate}
\item \textbf{Optimal $k$ scheduling:} Should different layers use different ranks?
\item \textbf{Interaction with LoRA placement:} Should LoRA be on down or up?
\item \textbf{Rank selection by singular value:} Can we auto-select $k$ per layer?
\item \textbf{Composition with KV compression:} How do they interact?

\end{enumerate}
